\subsection{Polarity of finite Brownian chord diagrams}
\label{arm:finite-diagram-subsection}

Fix a positive definite covariance matrix $\Sigma$ whose two coordinate
variances agree and whose correlation is $\rho\in(0,1)$.  Put
\[
 d_0=0,\qquad d_1=\frac12,\qquad
 d_2=\frac{\pi}{2\arccos(-\rho)},\qquad
 \beta=\frac12+d_2>1.
\]
The argument below uses only Brownian survival and landing estimates.
It does not use a lattice arm estimate, a random address code, or a
passage kernel.

Let $J_i=[a_i^-,a_i^+]$, $1\le i\le m$, be pairwise disjoint compact
intervals in increasing order, and let $K_i$ be nonempty compact
subintervals of their interiors.  Fix $s>0$ smaller than half the
distance of every $K_i$ from $\partial J_i$.  On each $J_i$ take
\[
 X_i(t)=A_i+B_i(t-a_i^-),
\]
where the $B_i$ are independent Brownian motions with covariance
$\Sigma$.  The vector $(A_1,\ldots,A_m)$ may have any law independent
of these Brownian motions.  For absolute-continuity comparisons one
may take the $A_i$ to be independent nonsingular Gaussian vectors.

A \emph{completion} of these pieces is a continuous function
$f:[a_1^-,a_m^+]\to\mathbb R^2$ which agrees with $X_i$ on every
$J_i$.  No probability law for the missing gaps is imposed.
Let $G_1,G_2$ be graphs on $\{1,\ldots,m\}$, with loops discarded,
and write
\[
 e=\operatorname{rank}(G_1)+\operatorname{rank}(G_2),
 \qquad \operatorname{rank}(G_c)=m-\#\operatorname{components}(G_c),
\]
where isolated vertices count as components.  The graph is
\emph{realized} if there are $t_i\in K_i$ and a completion $f$ such
that, for each edge $\{i,j\}$ of $G_c$,
\begin{equation}
\label{arm:finite-diagram-chord}
 f_c(t_i)=f_c(t_j)=\min_{[t_i\wedge t_j,t_i\vee t_j]}f_c.
\end{equation}
Denote this event by $\mathcal E_G$.
In particular, a chord tests every retained part of its span, including
any intermediate retained windows.  Requiring only equality of its
endpoint heights would define a different event.

\begin{lemma}[Finite-cell Brownian diagram bound]
\label{arm:finite-diagram-bound}
\label{arm:finite-diagram-polarity}
For each $\varepsilon>0$ there are $C_\varepsilon<\infty$ and
$\delta_0>0$, depending on the fixed windows, $m$, $\Sigma$, and
$\varepsilon$, such that for $0<\delta<\delta_0$,
\begin{equation}
\label{arm:finite-diagram-probability}
 \mathbb P(\mathcal E_G)
 \le C_\varepsilon\delta^{\beta e-m-\varepsilon}.
\end{equation}
More precisely, there are measurable events
$\mathcal F_{G,\delta}\supset\mathcal E_G$ which satisfy the same
bound.  The constants are uniform in the law and values of the starting
vectors $A_i$.  Consequently $\mathcal E_G$ has probability zero
whenever $\beta e>m$.
\end{lemma}

\begin{proof}
We first state the elementary Brownian estimates used in the proof.
For a Brownian block of fixed duration $s$ and $k\in\{1,2\}$ constant
lower walls, let $\tau_k$ be the first violation, in the constrained
coordinates.  For $0<h\le1$,
\begin{align}
 \sup_{z\in[0,h]^k}\mathbb P_z(\tau_k>s)
       &\le C_s h^{2d_k},\label{arm:finite-diagram-survival}\\
 \sup_{z\in[0,\infty)^k}
 \mathbb P_z(\tau_k>s,\ Z(s)\in[0,h]^k)
       &\le C_s h^{k+2d_k}.\label{arm:finite-diagram-landing}
\end{align}
Fixed multiples of $h$ change only $C_s$.
Both bounds follow from Lemma~\ref{arm:killed-landing}, after translating
the constant walls to zero and fixing the positive duration.
The estimates remain valid conditionally when the walls and starting
point are known before the block is exposed.  An initial point below
a tested wall has zero probability of satisfying the gate.

Choose $\vartheta\in(0,1/2)$ with $2\beta e\vartheta<\varepsilon$, and put
$\eta=\delta^{1/2-\vartheta}$.  Cover each $K_i$ by intervals of a
deterministic mesh of width $\delta$, retaining the finitely many mesh
cells which meet $K_i$.  For sufficiently small $\delta$ these cells,
and an $s$-interval on either side of each cell, lie in $J_i$.
Let $\mathcal H_\delta$ be the event that both coordinate oscillations
on every such cell are at most $\eta$.  The reflection principle and
a union bound give
\begin{equation}
\label{arm:finite-diagram-mesh-error}
 \mathbb P(\mathcal H_\delta^c)
 \le C\delta^{-1}\exp\{-c\eta^2/\delta\}
 =C\delta^{-1}\exp\{-c\delta^{-2\vartheta}\}.
\end{equation}
This bound is uniform in all starting vectors.

Fix one cell $[a_i,b_i]$ for each $i$.  Suppose the graph is realized
with $t_i\in[a_i,b_i]$, and suppose $\mathcal H_\delta$ holds.
For one coordinate, the scalar minimum-chord relation is transitive.
Indeed all endpoints in a connected component have the same height,
and the union of the spans along a path between two vertices covers
the interval between their times; every such span stays above that
height.  Thus, in each component of $G_c$, the chronologically
consecutive selected times themselves form minimum chords.  Use
these sorted chains, with one comparison per unit of monochrome rank.
Let $k_i^-$ and $k_i^+$ be the numbers of colors with a predecessor
and successor, respectively, at vertex $i$.  Then
\begin{equation}
\label{arm:finite-diagram-incidences}
 k_i^-,k_i^+\in\{0,1,2\},\qquad
 \sum_i k_i^- =\sum_i k_i^+=e.
\end{equation}
Write $h_{i,c}=X_{i,c}(a_i)$.

Use the left slot $[a_i-s,a_i]$ and the right slot $[b_i,b_i+s]$.
All these slots are disjoint.  If color $c$ has predecessor $j$ at
$i$, let $H_c$ be their common actual chord height.  Then
$|h_{j,c}-H_c|\le\eta$ and $|h_{i,c}-H_c|\le\eta$.
The entire left slot lies between the two actual chord endpoints,
so, with the wall $w_{i,c}^-=h_{j,c}-\eta$,
\begin{equation}
\label{arm:finite-diagram-left-gate}
 \inf_{[a_i-s,a_i]}X_{i,c}\ge w_{i,c}^-,\qquad
 0\le X_{i,c}(a_i)-w_{i,c}^-\le3\eta.
\end{equation}
If color $c$ has a successor at $i$, the entire right slot likewise
lies in that chord span.  With $w_{i,c}^+=h_{i,c}-\eta$, cell
oscillation and the chord minimum give
\begin{equation}
\label{arm:finite-diagram-right-gate}
 \inf_{[b_i,b_i+s]}X_{i,c}\ge w_{i,c}^+,\qquad
 0\le X_{i,c}(b_i)-w_{i,c}^+\le2\eta.
\end{equation}
These are necessary gates involving only fixed-slot path restrictions,
not a conditioning on the actual selected times or on
$\mathcal H_\delta$.

Condition on all starting vectors and expose the independent Brownian
pieces in the order $J_1,\ldots,J_m$, each in increasing time.
Every left wall in Equation~\eqref{arm:finite-diagram-left-gate}
uses a representative in an earlier window.  Every right wall in
Equation~\eqref{arm:finite-diagram-right-gate} uses the already
exposed left endpoint of its cell.  Their heights are therefore
known before the respective slots.  Applying
Equations~\eqref{arm:finite-diagram-survival}
and~\eqref{arm:finite-diagram-landing} and the tower property bounds
the probability of all gates for this cell tuple by
\[
 C_m\eta^{\sum_i(k_i^-+2d_{k_i^-}+2d_{k_i^+})}.
\]
The bound is uniform in the starting vectors, so it also holds before
conditioning on them.  Since
$d_k\ge k d_2/2$ for $k\in\{0,1,2\}$,
Equation~\eqref{arm:finite-diagram-incidences} shows that the exponent
is at least
\[
 e+2d_2e=2\beta e.
\]
There are $O(\delta^{-m})$ tuples of cells.  Let
$\mathcal F_{G,\delta}$ be the union of the gate events for these
tuples, together with $\mathcal H_\delta^c$.  All its constituents
are measurable, and the preceding deterministic implications give
$\mathcal E_G\subset\mathcal F_{G,\delta}$.  Consequently
\[
 \mathbb P(\mathcal F_{G,\delta})
 \le C\delta^{\beta e-m-2\beta e\vartheta}
       +C\delta^{-1}e^{-c\delta^{-2\vartheta}}.
\]
Equation~\eqref{arm:finite-diagram-probability} follows by the choice
of $\vartheta$; the case $e=0$ also follows directly from the trivial bound
one.  If $\beta e>m$, choose $\varepsilon\in(0,\beta e-m)$ and
send $\delta$ to zero.

For completeness the realization event can also be taken measurable
without selecting a completion.  Replace every missing gap by affine
interpolation in each coordinate.  If any completion satisfies the
chords, every gap traversed by a chord has both its endpoint values
above that chord height, so affine interpolation also satisfies the
chords.  Thus a completion exists exactly when this fixed interpolation
has a realizing tuple.  Minimum values on compact intervals vary
continuously with the path and with the interval endpoints.  Taking
the minimum of the total nonnegative chord discrepancies over the
compact set $\prod_i K_i$ makes realization a measurable event.
\end{proof}

\begin{corollary}[Transfer and a finite class bound]
\label{arm:finite-diagram-transfer}
The polarity assertion of Lemma~\ref{arm:finite-diagram-bound} holds
under every law absolutely continuous with respect to the joint law
of the pieces and their starting vectors described above.
Under a density bound $L\le K$, its measurable mesh-cover events obey
\[
 \mathbb Q(\mathcal F_{G,\delta}\cap\{L\le K\})
 \le K C_\varepsilon\delta^{\beta e-m-\varepsilon}.
\]
Mere absolute continuity transfers polarity, without asserting the
same polynomial bound for the mesh-cover events under the untruncated
law.

Suppose a collection of eligible time points has a scalar chord
relation in each color, and every fixed finite family of pairwise
separated enlarged windows about distinct eligible points has this
joint absolute-continuity property.  Assume also that every tested
finite chord diagram is realized by a continuous completion in a
fixed order of those windows, or one of finitely many such orders.
Then, almost surely, every class generated by the two chord relations
has at most
\begin{equation}
\label{arm:finite-class-bound}
 \left\lfloor\frac{\beta}{\beta-1}\right\rfloor
\end{equation}
distinct eligible points, provided each finite connected subdiagram of
that class uses only eligible points and can be localized in a countable
family of such windows.
\end{corollary}

\begin{proof}
Absolute continuity transfers each null event.  The quantitative
truncated assertion is simply integration of the density over
$\mathcal F_{G,\delta}$.  There is no assertion of a uniform density
bound without the displayed truncation.

For a connected union $G_1\cup G_2$ on $m$ vertices, take a spanning
tree of the union.  Its edges of either color form an independent set
in that monochrome graph, so $e\ge m-1$.  If
$m>\beta/(\beta-1)$, then
$\beta e\ge\beta(m-1)>m$, and the diagram is polar.
There are finitely many two-colored graphs on a fixed finite vertex
set, and finitely many window orders.  Countable buffered localization
therefore excludes all such connected diagrams simultaneously.
If a generated class had at least
$m_0=\lfloor\beta/(\beta-1)\rfloor+1$ eligible vertices, finite chains
connecting them would give a finite connected graph on at least
$m_0$ vertices.  A spanning tree can be reduced by removing leaves to
a connected subgraph on exactly $m_0$ vertices.  Its edges are still
actual chords, contradicting the preceding exclusion.  This proves
Equation~\eqref{arm:finite-class-bound}.
\end{proof}

\begin{remark}[Separated windows on several lines]
\label{arm:finite-diagram-multiple-lines}
For finitely many off-center time points on finitely many lines, first
localize each point in a compact rational window bounded away from any
excluded center and from the other selected points on its line.  Retain
larger windows with positive buffers.  The argument applies whenever
their joint pair-path restrictions, including relative starting-height
vectors, are absolutely continuous with respect to independent
Brownian pieces with full-dimensional starting-height density.
An unrecorded common additive vector is harmless: the chord event is
translation invariant, and one can add an independent Gaussian common
vector when comparing full starting-height laws.  The chronological
order must be fixed in the finite diagram or enumerated among its
finitely many possibilities.

The completion requirement is substantive.  It includes lower-wall
conditions throughout every retained portion of a chord span; mere
endpoint equalities or endpoint-only arms do not supply these tests.
The assertion also requires joint absolute continuity of the pieces
and offsets, rather than just Brownian marginal laws or a density for
the offsets in isolation.  No claim is made at a pinned center lacking
this joint absolute continuity, or about chains through such a center.
\end{remark}
